% ---------------------------------------------------------------------
\subsection{Descrizione Dettagliata Casi d'Uso (Template Cockburn)}
\label{sec:cockburn}
% ---------------------------------------------------------------------

Questa sezione fornisce la formalizzazione testuale strutturata di un caso d'uso significativo secondo il template di Alistair Cockburn. Il caso d'uso selezionatp è rappresentativo delle funzionalità core del sistema ed è stato scelto per dimostrare la completezza della specifica.

% ---------------------------------------------------------------------
\subsubsection{Caso d'Uso \#5: Segnalazione Issue}
\label{subsec:cockburn-segnala-issue}
% ---------------------------------------------------------------------

\begin{longtable}{|p{0.28\linewidth}|p{0.65\linewidth}|}
\hline
\textbf{Use Case \#5} & \textbf{Segnalazione Issue} \\ \hline
\endfirsthead

\hline
\textbf{Use Case \#5} & \textbf{Segnalazione Issue} \\ \hline
\endhead

\hline
\endfoot

\textbf{Goal in Context} & Un utente autenticato oppure un amministratore vogliono segnalare una nuova issue per tracciare un problema, richiedere una funzionalità, porre una domanda o segnalare un problema di documentazione. \\ \hline

\textbf{Preconditions} & \\ \hline

\textbf{Success End Conditions} & La issue viene creata con successo e aggiunta alla lista delle issue del progetto con stato "ToDo". L'utente può visualizzare la issue appena creata e riceve conferma dell'operazione. \\ \hline

\textbf{Failed End Conditions} & Nessun cambiamento viene effettuato al sistema. La issue non viene creata se i dati obbligatori sono incompleti o non validi. \\ \hline

\textbf{Primary Actor} & Utente Autenticato oppure Amministratore \\ \hline

\textbf{Trigger} & L'utente clicca sul pulsante "Crea Nuova Issue" o "Segnala Issue" dalla dashboard o dalla lista issues. \\ \hline

\end{longtable}

\newpage

\paragraph*{Main Scenario}

\begin{longtable}{|c|p{0.38\linewidth}|p{0.48\linewidth}|}
\hline
\textbf{Step} & \textbf{User Action} & \textbf{System Response} \\ \hline
\endfirsthead

\hline
\textbf{Step} & \textbf{User Action} & \textbf{System Response} \\ \hline
\endhead

\hline
\endfoot

1 & Clicca su "Segnala nuova issue" dalla dashboard & \\ \hline

2 & & Mostra il form di creazione issue vuoto (M1) \\ \hline

3 & Compila i campi obbligatori: 
\newline $\bullet$ Titolo,
\newline $\bullet$ Descrizione,
\newline $\bullet$ Tipo,
\newline Opzionalmente, sceglie priorità e/o aggiunge etichette e/o immagini
\newline Clicca sul pulsante "Crea Issue" & \\ \hline

4 & & Reindirizza alla pagina di vista Kanban (M6) e mostra il messaggio di successo "Issue creata con successo" (M2) \\ \hline

\end{longtable}

\newpage
\subsubsection*{Estensione A: Campo Titolo mancante}

\begin{longtable}{|c|p{0.38\linewidth}|p{0.48\linewidth}|}
\hline
\textbf{Step} & \textbf{User Action} & \textbf{System Response} \\ \hline
\endfirsthead

\hline
\textbf{Step} & \textbf{User Action} & \textbf{System Response} \\ \hline
\endhead

\hline
\endfoot

3.1a & Compila tutti i campi tranne il \textbf{Titolo} e clicca "Crea Issue" & \\ \hline

3.2a & & Mostra messaggio di errore (M3) \\ \hline

3.3a & Clicca "OK" o chiude il messaggio di errore & \\ \hline

3.4a & & Il form rimane visualizzato con i dati già inseriti preservati e ritorna allo step 3 dello scenario principale \\ \hline

\end{longtable}

\subsubsection*{Estensione B: Campo Descrizione mancante}

\begin{longtable}{|c|p{0.38\linewidth}|p{0.48\linewidth}|}
\hline
\textbf{Step} & \textbf{User Action} & \textbf{System Response} \\ \hline
\endfirsthead

\hline
\textbf{Step} & \textbf{User Action} & \textbf{System Response} \\ \hline
\endhead

\hline
\endfoot

3.1b & Compila tutti i campi tranne la \textbf{Descrizione} e clicca "Crea Issue" & \\ \hline

3.2b & & Mostra messaggio di errore (M4) \\ \hline

3.3b & Clicca "OK" o chiude il messaggio di errore & \\ \hline

3.4b & & Il form rimane visualizzato con i dati già inseriti preservati e ritorna allo step 3 dello scenario principale \\ \hline

\end{longtable}

\newpage
\subsubsection*{Estensione C: Campo Tipo mancante}

\begin{longtable}{|c|p{0.38\linewidth}|p{0.48\linewidth}|}
\hline
\textbf{Step} & \textbf{User Action} & \textbf{System Response} \\ \hline
\endfirsthead

\hline
\textbf{Step} & \textbf{User Action} & \textbf{System Response} \\ \hline
\endhead

\hline
\endfoot

3.1c & Compila tutti i campi tranne il \textbf{Tipo} e clicca "Crea Issue" & \\ \hline

3.2c & & Mostra messaggio di errore (M5) \\ \hline

3.3c & Clicca "OK" o chiude il messaggio di errore & \\ \hline

3.4c & & Il form rimane visualizzato con i dati già inseriti preservati e ritorna allo step 3 dello scenario principale \\ \hline

\end{longtable}

\subsubsection*{Estensione D: Annullamento operazione}

\begin{longtable}{|c|p{0.38\linewidth}|p{0.48\linewidth}|}
\hline
\textbf{Step} & \textbf{User Action} & \textbf{System Response} \\ \hline
\endfirsthead

\hline
\textbf{Step} & \textbf{User Action} & \textbf{System Response} \\ \hline
\endhead

\hline
\endfoot

3.1d & Durante la compilazione del form, clicca su "Annulla" oppure sulla "X" per chiudere & \\ \hline

3.2d & & Chiude il form senza salvare, reindirizza alla pagina precedente con vista Kanban (M6)\\ \hline

\end{longtable}

\newpage